\section*{अध्याय \ref{ch:b1:sinusoidal-impedance} --- ज्यावक्रीय स्थायी व्यवस्था और प्रतिबाधा}

\begin{solution}{exo:b1:sinusoidal-impedance:1}
$\underline U = 5\eu^{j\pi/3} = 2.5 + 4.33j$ (V)। $3\cos\omega t +
4\sin\omega t$: $3 + 4\eu^{-j\pi/2} = 3 - 4j$, मापांक $5$, कोणांक
$-\arctan(4/3) = \qty{-0.927}{rad}$: अतः $5\cos(\omega t - 0.927)$।
\end{solution}

\begin{solution}{exo:b1:sinusoidal-impedance:2}
$R$: दोनों पर \qty{100}{\ohm}। $L\omega$: \qty{31.4}{\ohm} (\qty{50}{Hz}),
\qty{628}{\ohm} (\qty{1}{kHz})। $1/C\omega$: \qty{318}{\ohm}, \qty{15.9}{\ohm}।
\qty{50}{Hz} पर श्रेणी $RL$: $\abs{\underline Z} = \sqrt{100^2 + 31.4^2} =
\qty{105}{\ohm}$, \omterm{def:b1:wave-propagation:signal}{कला} $\arctan(0.314) = 17^\circ$।
\end{solution}

\begin{solution}{exo:b1:sinusoidal-impedance:3}
$U = 230\sqrt2 = \qty{325}{V}$। $I_{\mathrm{rms}} = 2000/230 = \qty{8.7}{A}$,
शिखर \qty{12.3}{A}। $\langle Ri^2\rangle = R\langle i^2\rangle =
RI_{\mathrm{rms}}^2$, \omterm{def:b1:sinusoidal-impedance:rms}{प्रभावी मान} की परिभाषा से ही।
\end{solution}

\begin{solution}{exo:b1:sinusoidal-impedance:4}
$f_0 = 1/(2\pi\sqrt{LC}) = \qty{5.03}{kHz}$; $Q = \sqrt{L/C}/R = 3.16$;
$\Delta f = f_0/Q = \qty{1.6}{kHz}$; $I_{\max} = E/R = \qty{10}{mA}$।
\end{solution}

\begin{solution}{exo:b1:sinusoidal-impedance:5}
$1/\sqrt2$ $\omega = 1/RC$ पर: $f_c = 1/(2\pi RC) = \qty{1.0}{kHz}$।
\qty{1.0}{kHz} पर: अनुपात $0.71$, \omterm{def:b1:wave-propagation:signal}{कला} $-45^\circ$। \qty{10}{kHz} पर:
$RC\omega = 10$, अनुपात $1/\sqrt{101} = 0.10$, \omterm{def:b1:wave-propagation:signal}{कला} $-84^\circ$।
\end{solution}

\begin{solution}{exo:b1:sinusoidal-impedance:6}
$\underline Z = 100 + j(250 - 100) = 100 + 150j$: $Z = \qty{180}{\ohm}$,
$\arg = 56^\circ$। $I = 1.0/180 = \qty{5.6}{mA}$, जो स्रोत से
$56^\circ$ पिछड़ती है। $U_R = \qty{0.56}{V}$, $U_L = \qty{1.39}{V}$, $U_C = \qty{0.56}{V}$;
\omterm{def:b1:sinusoidal-impedance:complex}{कला-सदिश}: $U_R$ $I$ के अनुदिश, $U_L$ ऊपर, $U_C$ नीचे, परिणामी \qty{1.0}{V}।
$L\omega > 1/C\omega$: यानी अनुनाद के ऊपर (प्रेरकीय)।
\end{solution}

\begin{solution}{exo:b1:sinusoidal-impedance:7}
$f_0 = 1/(2\pi\sqrt{LC})$: \qty{1.59}{MHz} (\qty{50}{pF}) से \qty{0.50}{MHz}
(\qty{500}{pF}) तक। \qty{1.00}{MHz} पर: $C = 1/(4\pi^2f_0^2L) = \qty{127}{pF}$।
$\Delta f = \qty{10}{kHz}$। \qty{1.009}{MHz} पर केंद्र: $x = 1.009$,
$Q(x - 1/x) = 1.8$, अनुपात $1/\sqrt{1 + 3.2} = 0.49$। संधारित्र का \omterm{def:b1:dc-circuits:voltage}{विभवांतर}
$QE = \qty{1.0}{mV}$।
\end{solution}

\begin{solution}{exo:b1:sinusoidal-impedance:8}
$Q(x - 1/x) = \pm1 \Leftrightarrow x^2 \mp x/Q - 1 = 0$; धनात्मक मूल
$x_{1,2} = \mp 1/2Q + \sqrt{1 + 1/4Q^2}$। अंतर $1/Q$, अतः
$\Delta\omega = \omega_0/Q$; गुणनफल $(1 + 1/4Q^2) - 1/4Q^2 = 1$, अतः
$\omega_1\omega_2 = \omega_0^2$ (बैंड लघुगणकीय पैमाने पर सममित है)।
\end{solution}

\begin{solution}{exo:b1:sinusoidal-impedance:9}
$I = P/(U\cos\varphi) = 2000/(230 \times 0.70) = \qty{12.4}{A}$। \omterm{thm:b1:sinusoidal-impedance:power}{प्रतिघाती
शक्ति} $P\tan\varphi = 2000 \times 1.02 = \qty{2.0}{kvar}$। $C = 2040/(314
\times 230^2) = \qty{123}{\micro F}$। नई धारा $2000/230 = \qty{8.7}{A}$।
हानियाँ पहले $0.50 \times 12.4^2 = \qty{77}{W}$, बाद में $0.50 \times 8.7^2 =
\qty{38}{W}$।
\end{solution}

\begin{solution}{exo:b1:sinusoidal-impedance:10}
$\underline Y = 1/R + j(C\omega - 1/L\omega)$; $\underline U = \underline I/
\underline Y$ तब अधिकतम है जब प्रतिग्राहिता लुप्त हो, यानी $\omega_0$ पर,
जहाँ $U = RI$। अर्ध-शक्ति: $\abs{C\omega - 1/L\omega} = 1/R$, यानी
$RC\omega_0(x - 1/x) = \pm1$: अतः $\Delta\omega = \omega_0/(RC\omega_0) = 1/RC$
और $Q = RC\omega_0 = R\sqrt{C/L}$। आँकड़े: $\omega_0 = \qty{1.0e6}{rad/s}$
($f_0 = \qty{159}{kHz}$), $Q = 10^4\sqrt{10^{-6}} = 10$, $\Delta f =
\qty{16}{kHz}$।
\end{solution}

\begin{solution}{exo:b1:sinusoidal-impedance:11}
$U_C = I/C\omega = E/[C\omega\sqrt{R^2 + (L\omega - 1/C\omega)^2}]$;
भीतर $C\omega$ से गुणा कीजिए: $\sqrt{(RC\omega)^2 + (LC\omega^2 - 1)^2} =
\sqrt{x^2/Q^2 + (1 - x^2)^2}$। $s = x^2$ के साथ: $(1 - s)^2 + s/Q^2$
$s = 1 - 1/2Q^2$ पर न्यूनतम है (यदि वह धनात्मक हो), और न्यूनतम $1/Q^2 - 1/4Q^4$:
अतः $U_{C,\max} = QE/\sqrt{1 - 1/4Q^2}$। $Q = 3.16$: $x_r = 0.975$,
$U_{C,\max} = 3.2E$। $Q < 1/\sqrt2$ के लिए न्यूनतम $s = 0$ पर है: तब $U_C$
$E$ से एकदिष्ट रूप से घटता जाता है, कोई शिखर नहीं।
\end{solution}

\begin{solution}{exo:b1:sinusoidal-impedance:12}
अनुनाद पर $i = I\cos\omega_0 t$ और $u_C = (I/C\omega_0)\sin\omega_0 t$,
अतः $1/C\omega_0^2 = L$ का उपयोग करके,
\[
\tfrac12 Li^2 + \tfrac12 Cu_C^2 = \tfrac12 LI^2\cos^2\omega_0 t
+ \tfrac12 LI^2\sin^2\omega_0 t = \tfrac12 LI^2 .
\]
प्रति आवर्तकाल क्षय: $\tfrac12 RI^2T$। अनुपात $\times 2\pi$: $2\pi L/(RT)
= L\omega_0/R = Q$।
\end{solution}

\begin{solution}{pb:b1:sinusoidal-impedance:1}
\textbf{1.} $\underline Z = R + j(L\omega - 1/C\omega)$.

\textbf{2.} $\underline I = E/\underline Z$; $I = E/\sqrt{R^2 + (L\omega -
1/C\omega)^2}$; $\varphi_i = -\arctan[(L\omega - 1/C\omega)/R]$.

\textbf{3.} \omterm{def:b1:sinusoidal-impedance:impedance}{प्रतिघात} $\omega_0 = 1/\sqrt{LC}$ पर लुप्त हो जाता है: अतः $I_{\max}
= E/R$, $\varphi_i = 0$।

\textbf{4.} $C = 1/(4\pi^2f_0^2L) = \qty{127}{pF}$; $L\omega_0 = 2\times10^{-4}
\times 6.28\times10^6 = \qty{1257}{\ohm}$.

\textbf{5.} $Q = 1257/12.6 = 100$।

\textbf{6.} $L\omega - 1/C\omega = L\omega_0(x - 1/x)$, अब $R$ से भाग दीजिए।

\textbf{7.} $x < 1$: \omterm{def:b1:sinusoidal-impedance:impedance}{प्रतिघात} ऋणात्मक, $\varphi_i > 0$, धारा
आगे रहती है (धारिता जैसा); $x > 1$: पिछड़ती है (प्रेरकत्व जैसा)।

\textbf{8.} $Q(x - 1/x) = \pm1$.

\textbf{9.} $x^2 \mp x/Q - 1 = 0$, $x_{1,2} = \mp1/2Q + \sqrt{1 + 1/4Q^2}$,
$x_2 - x_1 = 1/Q$: अतः $\Delta\omega = \omega_0/Q$।

\textbf{10.} $\Delta f = \qty{10}{kHz}$.

\textbf{11.} \qty{1.010}{MHz}: $Q(x - 1/x) = 100 \times 0.0199 = 2.0$,
अनुपात $1/\sqrt{5.0} = 0.45$, शक्ति $0.20$: यानी \qty{-7}{dB}। \qty{1.020}{MHz}:
$3.96$, अनुपात $0.25$, शक्ति $0.06$: यानी \qty{-12}{dB}।

\textbf{12.} सीमांत: पड़ोसी केंद्र केवल \qty{7}{dB} नीचे है। ऊँचा
$Q$ (कम $R$), या शृंखला में कई सुर-बिठाए चरण, चुनाव को तीखा कर देते हैं।

\textbf{13.} \qty{0.50}{MHz} से \qty{1.59}{MHz} तक: यानी AM बैंड।

\textbf{14.} $\underline U_C = \underline I/jC\omega$; और $\omega_0$ पर
$U_C = I_{\max}/C\omega_0 = E/(RC\omega_0) = QE$।

\textbf{15.} $I_{\max} = 10^{-5}/12.6 = \qty{0.79}{\micro A}$; $U_C = QE =
\qty{1.0}{mV}$: यानी बिना किसी प्रवर्धक के सौ गुना विभव-लाभ।

\textbf{16.} $\underline U_L = jL\omega_0\underline I = jQE$, $\underline U_C
= -jQE$: बराबर और विपरीत, अतः वे कट जाते हैं; और स्रोत को अकेला $R$ दिखता है।

\textbf{17.} संचित $\tfrac12 LI_{\max}^2 = \qty{6.2e-17}{J}$; प्रति आवर्तकाल
क्षय $\tfrac12 RI_{\max}^2T = \qty{3.9e-18}{J}$; और $2\pi \times 6.2/0.39
\approx 100 = Q$।

\textbf{18.} $\tau = 2Q/\omega_0 = \qty{32}{\micro s}$। परिपथ
$\sim\tau$ से तेज़ बदलावों का पीछा नहीं कर सकता: बैंड-चौड़ाई $\Delta f$
केवल $\Delta f$ से धीमे मॉडुलनों को जाने देती है; $Q = 1000$ के साथ $\Delta f =
\qty{1}{kHz}$ होता और कार्यक्रम के ऊँचे स्वर खो जाते।

\textbf{19.} $Q$ आधा होकर $50$, $\Delta f$ दुगुना होकर \qty{20}{kHz}, और $U_C$
आधा होकर \qty{0.5}{mV}।

\textbf{20.} $R_{\text{कुल}} = \qty{62.6}{\ohm}$: $Q = 1257/62.6 = 20$,
$\Delta f = \qty{50}{kHz}$, $U_C = 20 \times \qty{10}{\micro V} = \qty{0.2}{mV}$:
यानी सुग्राहिता और वरणात्मकता, दोनों ढह जाती हैं।

\textbf{21.} $P = E^2/(2R_{\text{कुल}}) = 10^{-10}/125 = \qty{8e-13}{W}$;
और $R_a$ में अंश: $50/62.6 = 80\%$।

\textbf{22.} मिलान $R_a$ के बराबर भार-प्रतिरोध माँगता है, यानी कोई
बड़ा कुल प्रतिरोध; जबकि वरणात्मकता सबसे छोटा संभव कुल
प्रतिरोध माँगती है ($Q = L\omega_0/R_{\text{कुल}}$)। सीधे जोड़ से दोनों एक
साथ नहीं मिल सकते।

\textbf{23.} कोई टैप या कोई छोटा युग्मन संधारित्र ऐंटेना के सामने
अनुनादी परिपथ का केवल एक अंश रखता है: अतः ऐंटेना का प्रतिरोध
परिपथ को कहीं घटा हुआ ``दिखता'' है, जिससे $Q$ ऊँचा बना रहता है, जबकि
ऐंटेना फिर भी अपनी इष्टतम के पास शक्ति देता रहता है।

\textbf{24.} अर्ध-शक्ति आवृत्तियों पर \omterm{def:b1:sinusoidal-impedance:impedance}{प्रतिघात} $\pm R$ के बराबर होता है,
अतः $\varphi = \pm45^\circ$ और $\cos\varphi = 1/\sqrt2$; और साथ ही $I =
I_{\max}/\sqrt2$ भी, इसलिए $P = \tfrac12 EI\cos\varphi$ अपने अनुनादी मान का
आधा रह जाता है।

\textbf{25.} $Q$ संकेत को $Q$ गुना कर देता है और बैंड को
$f_0/Q$ तक सँकरा, जिसकी क़ीमत है धीमी अनुक्रिया ($\tau = 2Q/\omega_0$) और
ऐंटेना से नाज़ुक युग्मन; और \qty{1}{MHz} पर \qty{10}{kHz} चैनलों के लिए
$Q \approx 100$ ही वह संख्या है --- बशर्ते ऐंटेना हलके से
युग्मित रखा जाए।
\end{solution}
